\chapter{Fundamental Knowledge \& Related Work}
\section{Fundamental Knowledge}
This section discusses the fundamental techniques and applications foundational to the construction of the proposed system, categorized into 4 pillars: 
\begin{itemize}
    \item {Natural Language Processing (NLP):} The process of transforming raw, unstructured text into structured data formats that machines can process. This forms the basis for document processing and subsequent analysis.
    \item {Large Language Models (LLMs):} The reasoning components responsible for interpreting user queries, orchestrating tools, and generating responses. This enables the system to perform complex tasks and adapt to user needs.
    \item {Agentic Orchestration:} The management layer that coordinates specialized AI agents to execute complex, non-linear workflows.
    \item {System Design:} The architectural structure that ensures scalability and robustness, using design patterns tailored for educational environments.
\end{itemize}
\subsection{Natural Language Processing}

Natural Language Processing (NLP) converts unstructured document data into a structured representation. The goal is to translate raw text into a semantically precise format, enabling machine-inferable logic and further inferences.

The standard NLP pipeline follows this sequence:
\begin{center}
    {Information Extraction (IE) $\longrightarrow$ Knowledge Modeling $\longrightarrow$ Information Retrieval}
\end{center}

\subsubsection*{Information Extraction (IE)}
Information Extraction isolates relevant portions of text to produce a structured representation~\cite{StanfordNLP}. It filters out linguistic noise to capture the atomic units of knowledge. Traditional algorithms extract keywords using two primary methods:

\begin{itemize}
    \item {Keyword and Feature Extraction:} Uses dependency parsing and statistical methods (e.g., TF-IDF) to identify high-signal terms based on syntactic importance.
    \item {Named Entity Recognition (NER):} Classifies entities (e.g., academic theories, formulas, historical figures) within a text database~\cite{StanfordNLP}. These entities form the "nodes" of the system's knowledge graph.
\end{itemize}

To capture relationships that rigid rule-based systems miss, this process uses probabilistic modeling. It begins with Tokenization, breaking raw text into atomic units (tokens). Statistical algorithms like Conditional Random Fields (CRF) or Hidden Markov Models (HMM) then calculate probabilities for different interpretations, resolving linguistic ambiguities to predict the most likely semantic structure.

Traditional probabilistic models often struggle with sparse data and lack deep semantic understanding. Deep Learning addresses this by generating Embeddings. Instead of calculating discrete probabilities, these architectures map tokens into a continuous, high-dimensional vector space. In this space, geometrically close vectors represent semantically related concepts, allowing the system to capture context and nuance rather than relying solely on surface-level statistics.
\begin{figure}[H]
    \centering
    \includegraphics[width=0.7\textwidth]{Images/fundamental/embedding.png}
    \vspace{12pt}\caption{Transformers convert raw text into embeddings, which are numeric representations of semantic meaning}\label{fig:trans_embedding}
\end{figure}
The Transformer architecture has become the standard method in NLP, enabling modern Large Language Models (LLMs). Its core innovation is the Self-Attention mechanism. Rather than processing tokens sequentially, Transformers process the entire sequence simultaneously. The attention mechanism calculates a weighted importance score for every token relative to all other tokens, allowing the model to capture global context and long-range dependencies from unstructured texts.

\subsubsection*{Knowledge Modeling}
Often overlooked in standard Retrieval-Augmented Generation (RAG) systems, Knowledge Modeling is the construction stage where extracted entities are organized into a coherent structure that shapes how the system reasons.

\begin{itemize}
    \item {Semantic Parsing and Logical Form:} Maps extracted relations into a semantically precise form, such as First-Order Predicate Calculus (FOPC) or Knowledge Graphs (KG). This creates a deterministic "Backbone Graph" of the curriculum.
    \item {Embedding Space Topology:} Models knowledge within a continuous vector space where geometric proximity represents semantic similarity. This allows the system to handle concepts that are conceptually related but use different vocabulary.
\end{itemize}

\subsubsection*{Information Retrieval (IR)}
Information Retrieval is the active execution stage based on the results of IE and Knowledge Modeling. The goal is to identify data points relevant to a specific information need~\cite{StanfordNLP}.

\begin{itemize}
    \item {Lexical and Keyword Matching:} Uses exact term overlaps to find specific references or symbolic links within the knowledge base.
    \item {Vector Retrieval:} The standard modern approach, which uses the embedding space to retrieve context based on semantic similarity, typically calculated via Cosine Similarity:
    \begin{equation}
        \text{similarity}(\mathbf{A}, \mathbf{B}) = \frac{\mathbf{A} \cdot \mathbf{B}}{\|\mathbf{A}\| \|\mathbf{B}\|}
    \end{equation}
    \item {The GraphRAG Paradigm:} A retrieval method that navigates the entity-relationship graph established during the modeling phase. Unlike discrete chunk retrieval, GraphRAG performs "multi-hop" reasoning to connect disparate concepts, enabling the system to resolve complex queries that require a global understanding of the curriculum.
\end{itemize}


\subsection{Large Language Models (LLMs)}

Large Language Models (LLMs) are deep learning-based probabilistic models trained on massive datasets to understand and generate human-like text. Mathematically, an LLM performs the task of next-token prediction by modeling the conditional probability distribution of a token $w_{t}$ given a sequence of preceding tokens $(w_{1}, w_{2}, \dots, w_{t-1})$:

\begin{equation}
P(w_{t} | w_{1}, w_{2}, \dots, w_{t-1})
\end{equation}

The structural foundation of modern LLMs is the {Transformer} architecture, which relies on the {Self-Attention} mechanism. This mechanism allows the model to compute the relative importance of different tokens in a sequence regardless of their distance. The attention score is calculated using Query ($Q$), Key ($K$), and Value ($V$) matrices derived from input embeddings:

\begin{equation}
Attention(Q, K, V) = \text{softmax}\left(\frac{QK^T}{\sqrt{d_k}}\right)V
\end{equation}

By utilizing multi-head attention and positional encoding, Transformers can capture complex linguistic patterns and long-range dependencies more effectively than previous sequential architectures like Recurrent Neural Networks (RNNs).



[Image of Transformer architecture diagram]


\subsubsection*{LLM Ecosystem: Proprietary vs Open-Weights}

The landscape of Large Language Models in the industrial sector is divided by deployment strategy and architectural transparency. Closed-source, cloud-based models like GPT-4 and Claude form the proprietary tier. Provided as a Service (SaaS), these models are accessed via APIs hosted on centralized infrastructure. While they deliver state-of-the-art performance, their internal mechanisms (weights, datasets, and configurations) remain opaque. Additionally, per-token API costs can limit large-scale or iterative research applications.
\begin{figure}[H]
    \centering
    \includegraphics[width=0.7\textwidth]{Images/fundamental/cloud_local_llm.png}
    \vspace{12pt}\caption{Comparison of LLM Deployment Strategies: Traditional approaches (a) Direct cloud processing, (b) Local-only processing, and (c) Local-Cloud framework balance performance and security without leaking IP (proposed by~\cite{IP_Safe_llm})}\label{fig:cloud_local_llm}
\end{figure}

In contrast, open-weights models enable decentralized execution and architectural transparency. Projects like Meta's Llama series and Mistral release their pre-trained parameters, allowing deployment on private hardware. This ensures full control over data sovereignty and privacy, as information remains local. These models serve as baselines for academic research and industrial applications without relying on external APIs or subscription fees.

Building on these general-purpose baselines, researchers have fine-tuned models for specialized tasks. Coding specialists like Qwen-coder or CodeLlama are optimized for source code synthesis and logical reasoning, often rivaling larger general-purpose models in programming contexts. Beyond text, models have been extended into multimodal domains. By integrating vision encoders, developers enable visual comprehension and cross-modal reasoning.

\subsubsection*{LLM Inference}

This subsection explains the utility of LLM inference parameters in controlling the generation process. During inference, the next-token probability distribution predicted by the model can be dynamically altered to balance determinism and creativity.

\paragraph*{Temperature}
Temperature ($\tau$) controls the sharpness or randomness of the predicted probability distribution. Given the logits $z_i$ output by the final layer, the adjusted probability $P(w_i)$ for each candidate token is computed using the modified softmax function:
\begin{equation}
P(w_i) = \frac{\exp(z_i / \tau)}{\sum_{j} \exp(z_j / \tau)}
\end{equation}
A lower temperature ($\tau < 1$) sharpens the distribution, making the model more confident in its top choices and producing highly deterministic text, which is ideal for rigid tasks like JSON extraction. A higher temperature ($\tau > 1$) flattens the distribution, encouraging diversity and creativity.

\paragraph*{Top-k Sampling}
Top-$k$ sampling restricts the model from considering the entire vocabulary, instead keeping only the $k$ highest-probability candidates. This approach is conceptually similar to beam search, where only the most promising paths are kept active:
\begin{equation}
V_k = \mathop{\mathrm{arg\,max}}\limits_{V' \subset V, |V'|=k} \sum_{w \in V'} P(w)
\end{equation}
The probability mass is then redistributed among these $k$ tokens. This truncation prevents the model from generating highly improbable and potentially nonsensical words (the "long tail" of the distribution).

\paragraph*{Top-p (Nucleus) Sampling}
Top-$p$ (nucleus) sampling provides a dynamic alternative to Top-$k$. Instead of keeping a fixed number of tokens, it keeps the smallest prefix subset of the vocabulary whose cumulative probability mass meets or exceeds a threshold $p$:
\begin{equation}
\min \left\{ |V'| : V' \subset V, \sum_{w \in V'} P(w) \geq p \right\}
\end{equation}
This allows the candidate pool to expand or contract organically based on the model's confidence. When the model is very confident, the nucleus contains only a few tokens; when uncertain, the nucleus expands to include more options.

\paragraph*{Repetition Penalty}
Repetition penalty modifies the logits of tokens that have already appeared in the generated sequence, mathematically reducing their likelihood of being selected again. This prevents the model from entering infinite generation loops or excessively repeating specific phrases, ensuring a more natural text flow.


\subsection{Agentic Systems}

\subsubsection*{AI Agents}
An Artificial Intelligence (AI) Agent is a computational system capable of autonomous task execution through the dynamic design of workflows using available tools. Beyond language processing, AI agents can make decisions, solve problems, and interact with external environments.

\begin{figure}[H]
    \centering
    \includegraphics[width=0.7\textwidth]{Images/fundamental/agent.png}
    \vspace{12pt}\caption{Basic Single Agent diagram}\label{fig:react}
\end{figure}

Unlike standard LLMs, which rely solely on static training data and fixed knowledge limits, agentic technology uses tool calling on the backend to obtain up-to-date information. Agents optimize workflows and create subtasks autonomously, translating high-level natural language queries into executable workflows that complete real-life tasks. 

What truly pushes an AI agent to achieve what an LLM alone cannot is the harnessing of agentic technology. This involves constructing and integrating systems that boost capabilities, enabling the agent to navigate the complex, multi-step processes required for deep reasoning.

\paragraph*{LLM: The Logical Core}
Large Language Models (LLMs) serve as the brain of the agent, representing its thinking capacity. They provide the reasoning capabilities necessary to interpret vague user queries and orchestrate the appropriate tools. The LLM is responsible for parsing intent, identifying missing information, and deciding on the next actions. System decision-making depends heavily on the model used:
\begin{itemize}
    \item {Reasoning Models:} Models like OpenAI o1 or DeepSeek R1 are optimized for Chain of Thought processing. They excel at complex logical deduction and planning multi-step workflows, making them ideal for decomposing difficult queries into clear, executable plans.
    \item {Function Calling Models:} Models such as GPT-4o or Claude 3.5 Sonnet are fine-tuned to output structured data (e.g., JSON). This is efficient for tool configuration and defining output formats, which is critical for the execution phase.
    \item {Multimodal Models:} Models such as Gemini 1.5 Pro natively understand both text and visual inputs.
    \item {Lightweight Efficiency Models:} Smaller models like Llama 3 8B or GPT-4o-mini are faster and cheaper to run, making them ideal for narrow, repetitive tasks. 
\end{itemize}


\paragraph*{Context Engineering}
Context engineering represents the memorization capacity. It involves crafting the environment in which the LLM operates to ensure accurate and consistent outputs. This field includes several topics that govern how the model interacts with data and users:
\begin{itemize} 

\item {Prompting Strategies:}
Prompting is the primary method to instruct the agent, helping it understand the goals and constraints it is dealing with. Advanced methods like Chain of Thought (CoT) and Language of Thought (LoT) prompt the agent to expand its reasoning process step-by-step, ensuring a meaningful thought process before a final answer is generated.

\item {State Management:}
To maintain a consistent workflow, systems must store chat history and the current state. This allows the agent to understand the conversation context and the results from tools or other agents. State management ensures the agent knows exactly where it is in a multi-step workflow process, preventing redundant actions or confusion.

\item {Retrieval-Augmented Generation (RAG):}
Agents cannot remember all the information contained in massive data repositories, so RAG is used to retrieve external information that is not stored within the LLM's internal memory. RAG ensures the agent gets the needed information by querying relevant data segments or metadata to ground its answers in factual evidence.
\end{itemize}

\paragraph*{Tool Integration}
Representing the action capacity, tool integration is the process of connecting external resources (systems, software, hardware) or APIs to the agent, allowing it to perform specific functions that are beyond the capabilities of the LLM alone. This is critical for an automated system, as it enables the agent to truly interact with the external world, understand, and complete its mission from practice instead of static texts. However, it also poses significant threats of misuse when the agent makes real-life mistakes or fails to judge accurately.

\subsubsection*{Agent Frameworks}
The concept of an Agentic Framework arises from a practical necessity: building reliable autonomous systems requires more than just smart models; it requires a strong structure. Retraining Large Language Models to handle every specific rule or task is extremely expensive and inefficient. Therefore, the industry has shifted towards Agent Frameworks, which are software architectures designed to control and improve the behavior of AI without changing the model itself.

These frameworks act as a bridge, connecting the generative nature of AI with the reliable logic of software architecture. Agents are far more than LLMs with tools and context injection; they are intelligent systems with LLMs as the neural core and architecture following software logic. Modern frameworks leverage three main aspects in Agentic design and orchestration:

\paragraph*{{Multi-agent Orchestration}} 
In complex operations, giving a single agent too much information and too many different tasks often causes errors. A single agent has a limited attention span and context window. Frameworks solve this by using Multi-agent Systems. By breaking a large job into smaller parts handled by different agents, developers can assign specific agents to each task (e.g., one for visual analysis, one for audio processing, and one for synthesis), ensuring each agent stays focused and accurate.

\paragraph*{{Agent Memorization}}
Intelligent systems need to remember information over time, but standard LLMs treat every question as a new conversation (stateless). Frameworks solve this by building efficient memory systems with hierarchy, often using classes like Session or Memory. These classes act as data handlers containing multiple attributes and methods to help agents interact with memory to retrieve previous context during run-time. Memory management almost always follows these paradigms:
\begin{itemize}
    \item {Structured Memory:} Information is organized structurally, allowing the system to access relevant context precisely without overwhelming the agent with irrelevant details.
    \item {Context Window Optimization:} Frameworks implement strategies to optimize the context window. The memory is connected to an external database (usually MySQL), and agentic systems get needed context through lazy memory loading of tool results, conversation topics, etc. This performs dynamic pruning and memory injection, ensuring that the agent has access to the most relevant information without exceeding its limits, while also saving computational resources.
\end{itemize}

\paragraph*{{Tool Abstraction}}
Frameworks turn standard code functions into Tools that agents can use on their own. The key difference between a normal function and a Tool lies in its data structure: 
\begin{itemize}
    \item {Description:} Instructions on what the tool does and when to use it. The agent reads this description when needed without overwhelming itself with internal logic.
    \item {Input/Output Schema:} A defined structure for input and output parameters, often in JSON format, specifying the type and format of data the tool expects and returns. This allows the agent to call the tool correctly.
    \item {Logging memory:} Function results and errors are logged into memory for agent inference, context injection, and debugging.
\end{itemize}

Tool abstraction gives the AI capabilities like a calculator or a database search, allowing the agent to interact with the real world. Tool design must ensure tools are functional, easily interpretable by the agent, and have clear error handling to facilitate smooth interactions.

\subsubsection*{Langchain and LangGraph}\label{subsub:lang_intro}
The specific usecase of our educational system necessitates a flexible, nonlinear, and highly orchestrated approach. While numerous frameworks exist to harness the capabilities of Large Language Models (LLMs), the intricate requirements of educational guidance require a high-level, state-aware framework.

Unlike sequential or data-centric frameworks like LlamaIndex and SDK, LangGraph is built upon a \textit{State Graph Architecture} where the system is modeled as a directed graph. In this paradigm, each node represents a specialized agent or a discrete computational step, while the edges dictate the transition logic based on the system's current state. 

This granular control over the workflow enables non-linear branching, allowing the system to perform multipath reasoning for complex, non-repetitive tasks. In an educational context, because the workflow is globally state-aware, it can gracefully navigate diverse use cases without losing track of the student's long-term learning objectives. This is statistically difficult to achieve with stateless or purely sequential agentic frameworks.

\input{table/chap2/frameworks.tex}

The survey of various frameworks shown in Table~\ref{tab:frameworks} highlights LangGraph as the top framework for an educational system. While LlamaIndex offers stable operation with full data control, its data-centered nature often obscures the complex logic and asynchronous operations required for student-teacher interaction. The Google Generative AI SDK, though efficient for direct tool integration, lacks the complex state-management infrastructure necessary to handle unstable workflows compared to state-graph frameworks. The ability of LangGraph and CrewAI to maintain a persistent state and manage non-linear transitions provides the deterministic backbone required to facilitate an adaptive and reliable learning experience.
\subsection{System Design}

This section explains the software engineering principles, architectural patterns, and technological foundations used to construct scalable AI-integrated systems. 

Object-Oriented Programming (OOP) is a paradigm based on "objects" that contain data (attributes) and code (methods). In System Design, OOP is the industry standard because core architecture components are inherently stateful. Unlike simple functional scripts that execute and exit, system components must maintain persistent internal states, including conversation history and active tool configurations. OOP provides the natural structure to manage this state lifecycle. The system leverages the four main pillars of OOP:

\begin{itemize}
    \item {Abstraction:} Allows the system to model complex entities by hiding internal implementation details and exposing only relevant operations. This lets outside workflow logic interact with high-level methods rather than raw data from lower components.
    
    \item {Encapsulation:} Groups related variables and functions into a single unit, protecting the internal state from unintended interference.

    \item {Inheritance:} Promotes code reusability by allowing new classes to derive properties from existing ones. Specialized classes inherit these foundational capabilities while implementing unique behaviors, significantly reducing code duplication.

    \item {Polymorphism:} Allows objects of different classes to be treated as objects of a common superclass. In the orchestration layer, the system can iterate through a diverse list of components (like different retrieval engines or analysis modules) and trigger their execution loops using a uniform interface. The orchestrator does not need to know the specific component type, only that it adheres to the expected interface contract.
\end{itemize}

System architecture defines the overall structure of a software system, organizing its major components and how they communicate. Major applications often start with a Modular Monolithic Architecture. Unlike a traditional monolith where code is heavily tangled, a modular monolith compiles the application as a single unified program but keeps clear logical boundaries between different functions. This approach offers practical early-stage benefits: it simplifies deployment, ensures fast internal communication, and makes the codebase easier to test and debug.
\begin{figure}[H]
    \centering
    \includegraphics[width=0.7\textwidth]{Images/fundamental/micro.png}
    \vspace{12pt}\caption{Comparison of Microservice and Monolithic architecture}\label{fig:micro_vs_mono}
\end{figure}

However, as system demands grow, combining heavy processing tasks with basic web routing in a single program can cause performance bottlenecks. Microservices take the next step by turning logical modules into physical, independent services. Each service runs in its own process, often with its own separate database, and communicates with other services using standard protocols like REST APIs or message brokers. This architecture inherits the organizational advantages of the modular monolith while providing several additional merits:
\begin{itemize}
    \item {Independent Scalability:} Microservices allow different parts of the system to scale separately. For example, a heavy processing service can be upgraded with more computing power or duplicated without needing to scale the lightweight web server, ensuring efficient hardware use during high traffic.
    
    \item {Fault Isolation:} Software failures are contained. If a heavy processing service crashes due to memory limits, it does not bring down the entire system. Other services, such as user login and basic web browsing, remain fully operational while the crashed service restarts.
    
    \item {Technology Flexibility:} Microservices allow developers to choose the best programming language or tool for each specific task. For instance, data science components can be built using Python for its machine learning libraries, while high-speed API routing can be written in Go or Node.js.
\end{itemize}

Alongside these benefits, microservices possess some disadvantages. Internal function calls in a monolith are extremely fast. In microservices, these are replaced by network requests over the internet or local network, adding delay and processing overhead. Moreover, managing, monitoring, and deploying multiple independent services with separate databases is significantly harder than deploying a single program. It also requires complex strategies to keep data consistent across independent services.


\subsection{Technology Stack and Frameworks}
The selection of frameworks is driven by the need for high-performance systems with seamless AI integration. The following frameworks are effective building blocks for implementing the application layer.

\subsubsection*{Backend Development}
Backend development serves as the foundational logic of any software application, responsible for data processing, security, and the orchestration of communication between the database and the client. For an intelligent educational system involving heavy AI workloads, the backend must be capable of handling high-concurrency requests without latency, ensuring that complex agent computations do not block the user experience.
\begin{figure}[H]
    \centering
    \includegraphics[width=0.4\textwidth]{Images/system/backend/fastapi.png}
\end{figure}

To address these requirements, we utilize FastAPI, a modern web framework for building APIs with Python. While older frameworks like Django were built on synchronous standards known as WSGI, FastAPI adopts the Asynchronous Server Gateway Interface (ASGI). This approach allows the server to handle asynchronous code natively using Python's await syntax. This represents a paradigm shift for AI applications, moving away from blocking, sequential request handling to a concurrent model where I/O-bound tasks (such as waiting for an LLM response or a database query) release the CPU to handle other incoming requests, offering key features:
\begin{itemize}
    \item {High Concurrency:} By utilizing the ASGI standard, FastAPI can handle thousands of simultaneous connections (such as multiple users uploading videos or chatting with agents) without the blocking issues inherent in traditional synchronous frameworks.
    \item {Data Validation:} It integrates tightly with Pydantic, a data validation library. This ensures that the complex JSON outputs generated by AI agents are rigorously validated against strict schemas before being sent to the frontend, preventing runtime type errors.
    \item {Performance:} Built on top of Starlette, FastAPI offers execution speeds comparable to compiled languages like Go, which is critical for minimizing the latency overhead in the educational content processing pipeline.
\end{itemize}

\subsubsection*{UV: High-Performance Package Management}
Traditional Python dependency management with pip and virtualenv introduces significant overhead in complex AI projects, where dependency resolution across dozens of GPU-accelerated libraries can take several minutes. The system adopts {uv}, a Rust-based Python package manager developed by Astral, which consolidates virtual environment creation, dependency resolution, and lockfile generation into a single binary.
\begin{figure}[H]
    \centering
    \includegraphics[width=0.4\textwidth]{Images/system/backend/uv.png}
\end{figure}

UV achieves dependency resolution 10 to 100 times faster than pip by implementing a custom resolver in Rust. Its deterministic lockfile mechanism ensures that all team members and deployment environments operate with identical dependency versions, eliminating the "works on my machine" class of bugs. For this project, uv manages all backend dependencies, including LangChain, LangGraph, PyTorch, and the database client libraries, providing a consistent and reproducible build pipeline from development through production deployment.

\subsubsection*{Ollama: Local Model Inference and Optimization Infrastructure}
Ollama is an advanced abstraction layer designed for managing and executing LLMs on local infrastructure. It is built upon the \texttt{llama.cpp} core, providing an efficient environment for model orchestration and hardware acceleration.
\begin{figure}[H]
    \centering
    \includegraphics[width=0.4\textwidth]{Images/system/AI/ollama.png}
\end{figure}

A critical feature of local inference is {Quantization}, the process of reducing the precision of model weights from high-bit floating-point formats (e.g., FP16 or BF16) to lower-bit integers (e.g., 4-bit or 8-bit). This technique significantly reduces the VRAM footprint and memory bandwidth requirements:
\begin{equation}
w_{q} = \text{round}\left(\frac{w}{\Delta}\right)
\end{equation}
where $w$ is the original weight and $\Delta$ is the scaling factor. Quantization enables large-scale models to run on consumer-grade hardware with minimal degradation in perplexity and reasoning capabilities.

Moreover, it utilizes the {GGUF (GPT-Generated Unified Format)}, a binary format designed for rapid model loading and efficient metadata storage. Key technical characteristics include:
\begin{itemize}
    \item {mmap Compatibility:} Allows models to be mapped directly into memory, enabling nearly instantaneous startup times.
    \item {Unified Execution:} Supports flexible memory offloading, where layers of the model can be distributed between the GPU (VRAM) and CPU (System RAM) based on available resources.
    \item {Modelfile Abstraction:} Provides a structured way to define system prompts, temperature settings, and context window sizes within a single configuration file, ensuring deterministic behavior across different local environments.
\end{itemize}

\subsubsection*{The Langchain Ecosystem}
As mentioned in Subsection~\ref{subsub:lang_intro}, Langchain~\cite{LangChain} is a powerful framework for our educational system, with State Graph Architecture. Moreover, it is one of the most popular agentic frameworks in the industry, with a large community and rich documentation. The accumulated contributions to the ecosystem and agentic research using Langchain have provided us with a wealth of tools and integrations:
\begin{itemize}
    \item {Pre-built Agents:} Langchain offers a variety of pre-built agents that can be easily customized for specific use cases, such as question-answering, summarization, and multi-step reasoning. In 2026, it published DeepAgent, which is capable of extremely long-term reasoning and complex workflow orchestration.
    \item {Langfuse:} Langfuse, along with LangSmith, is one of the most powerful tools for agentic system development. It provides a comprehensive observability platform for monitoring and debugging agent interactions, allowing developers to trace the decision-making process of agents in real-time. This is critical for diagnosing issues in complex workflows and ensuring that the system behaves as expected, making agentic operation more transparent, manageable, and less error-prone. 
    \item {Context Management:} Langchain includes built-in context management capabilities, enabling agents to maintain context across interactions and manage state effectively. These include utilities for memory management (adapting to multiple external databases), conversation history, and dynamic context injection, which are essential for building complex agentic systems for long interactions over time.
\end{itemize}




\subsection{Database System Design \& Management}
A database is a structured system where applications store, read, write, update, and delete necessary information. In the domain of modern software engineering, data handling has become increasingly complex. To address this, robust architectures utilize a Polyglot Persistence strategy, which leverages distributed storage technologies optimized for different data modalities.

\begin{table}[H]
    \centering
    \caption{Polyglot Persistence Architecture: Strategic allocation of storage engines.}\label{tab:polyglot_db}
    \renewcommand{\arraystretch}{1.5}
    \small 
    \begin{tabularx}{\textwidth}{|p{2.0cm}|p{2.0cm}|X|X|}
    \hline
    {Storage Engine} & {Data Modality} & {System Role} & {Theoretical Justification} \\
    \hline
    Relational (MySQL) & Structured Data & User Authentication, Billing, Access Logs & Enforces ACID properties for transactional integrity and supports complex relational joins for administrative analytics. \\
    \hline
    NoSQL (MongoDB) & Semi-Structured Data & Semantic Metadata, Analysis Tags, JSON Outputs & Provides schema polymorphism to handle variable outputs from different AI agents (e.g., OCR vs. Object Detection) without sparse tables. \\
    \hline
    Object Storage (MinIO) & Unstructured Binary & Lecture Materials, Educational Assets, Audio Clips & Decouples storage from compute to prevent database bloat; utilizes Erasure Coding for high-throughput parallel streaming during ingestion. \\
    \hline
    Vector Database (Milvus) & High-Dimensional Vectors & Embeddings Index, Semantic Search & Optimized for Approximate Nearest Neighbor (ANN) algorithms (e.g., HNSW), enabling efficient similarity search in latent spaces where traditional B-Tree indexing fails. \\
    \hline
    Graph Database (Neo4j) & Multi-hop Graph Data & Knowledge Graph, Prerequisite Relationships & Optimized for traversing complex relationship networks; enables recursive queries and multi-hop reasoning to connect disparate educational concepts. \\
    \hline
    \end{tabularx}
\end{table}

\subsubsection*{Relational Database (RDBMS)}
Relational Databases are mathematically founded on Set Theory and First-Order Predicate Logic. They organize data into tuples (rows) and attributes (columns) within strictly defined relations (tables). The defining characteristic of this model is the enforcement of ACID properties (Atomicity, Consistency, Isolation, Durability), which guarantees that database transactions are processed reliably.

RDBMS is typically utilized for data that requires high integrity and deterministic relationships, such as User Management and Billing. It acts as the source of truth for entity relationships, enabling complex SQL JOIN operations (e.g., selecting all learning materials uploaded by a specific role within a timeframe). The adoption of RDBMS for core administrative data offers specific advantages to system stability.
\begin{figure}[H]
    \centering
    \includegraphics[width=0.4\textwidth]{Images/system/database/sqlite.png}
\end{figure}

MySQL is utilized as the primary Relational Database Management System. It handles core entities such as user accounts, roles, and administrative logs, ensuring that all structured interactions follow strict relational integrity.
\begin{itemize}
    \item {Transactional Integrity:} In a multi-user environment where credits or permissions are updated frequently, RDBMS ensures that race conditions do not occur. For instance, if a user purchases a subscription while simultaneously uploading an educational asset, the ACID properties ensure the account balance is updated accurately before the upload permission is granted.
    \item {Data Normalization:} By structuring data into normalized tables (e.g., Users, Roles, Logs), redundancy is minimized. This ensures that a change in a user's details is reflected system-wide instantly without needing to update thousands of duplicate records.
    \item {Complex Querying:} The robust SQL engine allows administrators to perform audit trails and analytics efficiently, such as tracking system usage patterns across different user groups, which is computationally expensive in non-relational systems.
\end{itemize}

\subsubsection*{Non-Relational Database (NoSQL)}
To address the limitations of rigid RDBMS schemas, Non-Relational databases emerged, often adhering to the CAP Theorem (Consistency, Availability, Partition Tolerance). Unlike RDBMS, which prioritizes strong Consistency, NoSQL databases like MongoDB often prioritize Partition Tolerance and Availability, offering a flexible schema-less design.

Document stores are typically the primary engine for storing semantic metadata derived from educational content analysis. Metadata derived from curriculum analysis is highly variable; one resource might contain OCR text, while another contains only Audio Transcripts or Face Recognition data. Document stores handle these as dynamic BSON (Binary JSON) documents, allowing fields to be added on the fly without the sparse matrices typical of rigid SQL tables.
\begin{figure}[H]
    \centering
    \includegraphics[width=0.4\textwidth]{Images/system/database/mongodb.png}
\end{figure}

MongoDB serves as the high-performance document store for managing the semi-structured metadata generated by AI agents. It allows the system to store diverse analysis results—such as transcripts, OCR results, and agentic reasoning logs—within a single, flexible schema.
\begin{itemize}
    \item {Schema Agnosticism:} AI agents are constantly evolving. If an "Object Detection" agent is upgraded to output new attributes (e.g., "Confidence Score" or "Bounding Box Coordinates"), document stores allow the new data to be saved immediately without migrating the entire database schema or breaking existing application logic.
    \item {Read Performance for Aggregates:} In a Knowledge-QA scenario, the system needs to retrieve the Video Title, Description, Transcripts, and Vector IDs simultaneously. In an RDBMS, this might require joining multiple tables. In a NoSQL system, this data can be stored as a single nested document, allowing the application to retrieve the entire context of a lesson in a single, high-speed read operation.
    \item {Horizontal Scalability:} As archives grow to thousands of hours of lecture content, metadata storage requirements increase. Native sharding capability allows data to be distributed across multiple servers economically, ensuring lookup times remain low even as the dataset grows.
\end{itemize}

\subsubsection*{High-Performance Object Storage}
While RDBMS and NoSQL excel at storing alphanumeric data, they are fundamentally inefficient for storing massive binary large objects (BLOBs) like lecture recordings, leading to database bloat and performance degradation. Object Storage systems like MinIO treat data not as a file hierarchy or a table, but as discrete objects stored in a flat address space.

High-performance architectures utilize Object Storage to handle the high-volume nature of raw educational media. It utilizes Erasure Coding—a mathematical algorithm that splits data into fragments and expands them with redundant data pieces—allowing for data recovery even if multiple drives fail. This approach decouples storage from compute; the heavy educational artifacts are offloaded to Object Storage, while the databases store only the lightweight reference pointers (URLs) to these objects.
\begin{figure}[H]
    \centering
    \includegraphics[width=0.3\textwidth]{Images/system/filestorage/minio.jpg}
\end{figure}

MinIO is deployed as a high-performance, S3-compatible object storage server. It is responsible for hosting all raw multimedia assets, including lecture videos, extracted images, and audio segments used for AI training and inference.
\begin{itemize}
    \item {Decoupling Storage and Compute:} By separating the raw files from the metadata databases, "database bloat" is prevented. This ensures that the RDBMS remains lightweight and responsive for user queries, while the Object Store handles the heavy I/O load of content streaming.
    \item {Parallel Streaming for AI Ingestion:} When an agent analyzes a lecture recording, it requires high-speed access to the file to extract frames. Object Storage supports parallelized byte-range requests, allowing the AI to read multiple parts of a media file simultaneously. This significantly reduces the time required for the ingestion and preprocessing phase compared to reading from a standard file server.
    \item {Cost-Efficiency and Redundancy:} Erasure Coding allows the system to survive hardware failures without data loss, while consuming less storage space than traditional replication methods. This is vital for maintaining a massive archive of media files economically.
\end{itemize}

\subsubsection*{Vector Database}
Traditional databases are designed for structured data and are inefficient for handling high-dimensional vector data generated by AI models. Vector databases are optimized for storing and querying high-dimensional vectors, which represent the semantic content of lecture segments. They utilize Approximate Nearest Neighbor (ANN) algorithms, such as Hierarchical Navigable Small World (HNSW) graphs, to enable efficient similarity search in latent spaces where traditional indexing methods fail.

The Vector Database is the backbone of the Retrieval-Augmented Generation (RAG) paradigm, allowing the system to perform semantic search based on the embeddings generated by the AI agents. When a user query is processed, it is transformed into an embedding vector, which is then compared against the vectors stored in the database to retrieve the most semantically relevant knowledge segments. The adoption of a Vector Database is critical for enabling the advanced retrieval capabilities required for educational analysis.

Milvus is the chosen vector database engine for SmartEdu. It indexes high-dimensional embeddings generated from course documents and lecture transcripts, providing sub-second retrieval times for semantic similarity queries across millions of vectors.
\begin{itemize}
    \item {Efficient Similarity Search:} Traditional databases cannot efficiently query high-dimensional vectors. Vector databases provide specialized indexing structures to handle this efficiently.
    \item {Scalability for Large Datasets:} As the number of educational segments grows, the vector database can scale horizontally to maintain low latency for similarity searches, ensuring that retrieval times remain fast even as the dataset expands.
    \item {Integration with AI Workflows:} Vector databases are designed to integrate seamlessly with AI pipelines, supporting real-time insertion and retrieval of embeddings.
    \item {Support for Dynamic Data:} As new educational materials are ingested and analyzed, their embeddings can be added to the vector database in real-time, allowing the system to continuously evolve and improve its retrieval capabilities without downtime.
\end{itemize}

\subsubsection*{Graph Database}
Neo4j is utilized to store the "Backbone Graph" of the knowledge base. It represents educational entities (concepts, lessons, modules) as nodes and their pedagogical relationships as directed edges. 
\begin{itemize}
    \item {Multi-hop Reasoning:} Unlike relational databases that require expensive JOIN operations, Neo4j allows for efficient traversal of deep relationship chains. This is critical for detecting knowledge gaps that may span multiple prerequisite levels.
    \item {Cypher Query Language:} The use of Cypher enables the system to express complex pedagogical patterns, such as "Find all prerequisite concepts of concept X that the student has not yet mastered," in a readable and performant manner.
    \item {Dynamic Schema:} The graph model allows the knowledge structure to evolve as new educational materials are added, without needing to redefine rigid table structures.
\end{itemize}

\subsubsection*{Database Management and Observation}
Managing multiple database engines requires a unified observation layer to ensure data consistency and facilitate debugging during development.
DbGate is an open-source, multi-database manager used as the primary dashboard for the SmartEdu infrastructure. It provides a unified web interface to inspect and manipulate data across MySQL and MongoDB instances simultaneously.
\begin{itemize}
    \item {Unified Interface:} Reduces the overhead of switching between multiple specialized management tools.
    \item {Schema Visualizer:} Helps developers understand the evolving structure of NoSQL documents and Relational tables in real-time.
    \item {Query Benchmarking:} Allows for testing and optimizing complex queries before they are integrated into the application's repository layer.
\end{itemize}

\subsection{Deployment and Containerization}
Deployment is the process of making a software application available for real-world usage. This section discusses the deployment strategy and the use of containerization to ensure that the system can be reliably and efficiently deployed for the AI pipeline and high-performance stacks discussed above.

Docker~\cite{docker_docs} is a containerization platform that enables developers to bundle applications and their required components into isolated units called {containers}. Each container acts as a lightweight, virtualized environment that includes all the necessary elements for the application to run. Because containers are portable and consistent, they work reliably across various environments, such as on-premise systems, cloud platforms, and clustered infrastructures.
\begin{itemize}
    \item {Portability:} Containers package the application together with all its dependencies, making it simple to move workloads between different environments without compatibility issues. Whether running on a local machine or a cloud server, a Docker container behaves identically, ensuring consistency throughout the development and deployment pipeline.
    \item {Isolation:} Docker provides strong isolation between applications and their dependencies. Each container operates independently, preventing interference between applications running on the same host and enhancing security and stability.
    \item {Resource Efficiency:} Containers are more resource-efficient than traditional virtual machines. Since containers share the same operating system kernel, they consume significantly less memory and disk space, reducing hardware overhead and simplifying resource allocation.
    \item {Scalability:} Docker supports efficient scaling of applications. By using orchestration tools such as Kubernetes or Docker Swarm, developers can easily launch multiple container instances to rapidly adjust capacity and meet varying levels of demand.
    \item {Faster Development and Deployment:} Developers can work in isolated, reproducible environments identical to production, minimizing integration issues. The deployment process becomes automated, streamlined, and less error-prone.
\end{itemize}

\subsection{Security: OAuth 2.0 Authorization}
OAuth is a widely used authentication and authorization protocol in backend systems across various programming languages and frameworks. The primary benefits of incorporating OAuth into backend systems include:
\begin{itemize}
    \item {Security:} OAuth provides a secure way for users to grant third-party applications limited access to their resources without sharing credentials, reducing the risk of unauthorized access and protecting user data.
    \item {User Experience:} It enables seamless authentication and authorization by allowing users to log in using existing accounts from trusted providers, eliminating the need to create and manage separate credentials.
    \item {Single Sign-On (SSO):} OAuth facilitates SSO capabilities, enabling users to access multiple applications and services with a single set of credentials and reducing password fatigue.
    \item {Granular Permissions:} It allows defining and enforcing granular permissions for accessing resources, ensuring that third-party applications only have access to specifically required resources.
    \item {Scalability:} By offloading user authentication and authorization to external providers, the backend system can better handle scalability challenges, as OAuth providers are equipped to process large volumes of authentication requests efficiently.
    \item {Standardization:} OAuth is a widely adopted standard supported by major technology companies and libraries, ensuring compatibility with a broad ecosystem of tools and services.
    \item {Compliance:} Implementations often include built-in support for regulatory requirements such as GDPR (General Data Protection Regulation) or HIPAA (Health Insurance Portability and Accountability Act), simplifying compliance efforts.
    \item {Developer Ecosystem:} OAuth has an active developer community providing extensive documentation, libraries, and support (such as PyJWT for protocol schemas, etc.), simplifying implementation and troubleshooting.
\end{itemize}

\section{Related Work}

\subsection{Document-Knowledge Modeling}
Converting documents into structured knowledge representations has long been a critical task in both industry and educational domains. To replace manual curation in educational contexts, which is time-consuming and not scalable, prior attempts have focused on entity extraction, prerequisite relation reasoning, and Retrieval-Augmented Generation (RAG). This has proven to be an effective strategy, demonstrated by early supervised models like \texttt{PREREQ}~\cite{PREREQ}, which utilized latent topic representations with Siamese networks to infer course dependencies and achieved over 70\% Precision@100 on the \texttt{University Course} and \texttt{NPTEL MOOC} datasets. However, traditional methods of entity extraction and relation reasoning often rely on rule-based systems or pure statistical machine learning techniques that struggle to capture the complexity and nuances of educational content relationships. Moreover, recent research~\cite{weller2025theoretical,LostInTheMiddle} demonstrates that models struggle to access relevant information when documents exceed a certain length. Consequently, although prior studies achieved impressive results on medium-scale retrieval benchmarks, they are difficult to apply to massive and unstructured enterprise datasets.
\begin{figure}[H]
    \centering
    \includegraphics[width=0.8\textwidth]{Images/relate/kggen.png}
    \vspace{12pt}\caption{KGGen with multiturn LLM reasoning, grouping generation}\label{fig:kggen}
\end{figure}
To address these limitations, LLM-based frameworks have been widely proposed. A prime example is \texttt{KGGen}~\cite{KGGen}, an automated system that utilizes a clustering algorithm to group synonymous entities and relations, resulting in denser and better-connected graphs that scored an 81.03\% average fact accuracy on the \texttt{MINE-2} (Measure of Information in Nodes and Edges) benchmark. The \texttt{EDC}~\cite{EDC} framework extends this capability by extracting, automatically defining schemas, and canonicalizing knowledge triples without predefined schemas, consistently surpassing state-of-the-art baselines on complex datasets like \texttt{WebNLG}, \texttt{REBEL}, and \texttt{Wiki-NRE}. 
\begin{figure}[H]
    \centering
    \includegraphics[width=0.8\textwidth]{Images/relate/edukg.png}
    \vspace{12pt}\caption{The end-to-end EduKGPipeline framework for educational knowledge construction and retrieval.}\label{fig:edukgpipeline}
\end{figure}
However, LLMs have significant drawbacks. First and most popular of all is the "hallucinations problem", leading to inaccurate and destructured extractions. To mitigate this, \texttt{GraphJudge}~\cite{GraphJudge} proposes using a fine-tuned open-source LLM as a "judge" to score and filter out incorrect knowledge triples generated by naive LLMs. Validated on general (\texttt{REBEL-Sub}, \texttt{GenWiki}) and domain-specific benchmarks (\texttt{SCIERC}, \texttt{Re-DocRED}), it achieved over 90\% judgment accuracy. Furthermore, for large-scale educational and enterprise systems requiring cost and latency optimization, comprehensive hybrid pipelines have proven highly effective in practice. For instance, \texttt{EduKGPipeline}~\cite{EduKGPipe} establishes an end-to-end framework fusing traditional text extraction with transformer-based weighting, demonstrating high precision on educational transcripts from the \texttt{CCI dataset}. 
\begin{figure}[H]
    \centering
    \includegraphics[width=0.8\textwidth]{Images/relate/graphrag hybrid.png}
    \vspace{12pt}\caption{Hybrid architecture of GraphRAG allow minimal cost and guarantee alignment of generated to original context}\label{fig:graph_hybrid}
\end{figure}
Similarly, efficient GraphRAG architectures~\cite{GraphRAG} substitute expensive LLM extractions with classical dependency parsing to maintain scalability, retaining 94\% of LLM-based performance on \texttt{RAGAS} metrics when evaluated on real-world enterprise datasets like \texttt{CCM} (Custom Code Migration). Finally, \texttt{CoDe-KG}~\cite{CoDeKG} utilizes syntactic decomposition to process complex sentences, thereby significantly enhancing information recall by an 8-point gain on rare relations across the \texttt{REBEL} and \texttt{WebNLG2} benchmarks.

Another inevitable error lies in the LLMs' limited context length, making them unable to generalize broad knowledge. They perform well when discovering short-dependent context, entities, or relationships in small sections, but fail for long dependencies in extremely long documents and course information. Trying to inject an excessive amount of information even when cut by windows, until now with the current level, only leads to redundant costs from multiple calls, information overload, and decreased accuracy. To uncover hidden relationships and deduce learning paths beyond explicit textual representations, several studies have shifted toward graph-algorithmic approaches.
\begin{figure}[H]
    \centering
    \includegraphics[width=0.8\textwidth]{Images/relate/dkg.png}
    \vspace{12pt}\caption{Dynamic Knowledge Graph (DKG) approach tracking knowledge evolution over time.}\label{fig:dkg}
\end{figure}
The \texttt{Dynamic Knowledge Graph (DKG)}~\cite{DKG} is a prominent representative of this approach, analyzing macroscopic structures within Knowledge Graphs. To recommend robust relational paths, DKG introduces the \texttt{KC2 (Knowledge Connector Candidate)} algorithm, which uses closeness centrality to suggest connections between disconnected knowledge islands. Additionally, to cluster knowledge into logical chapters or modules, the \texttt{Leiden}~\cite{Leiden} algorithm is utilized to optimize the Constant Potts Model (CPM) quality function, ensuring that knowledge communities remain well-connected rather than disjointed, addressing a common limitation of the traditional Louvain algorithm. It also divides graphs into discrete time snapshots and applies network science to track knowledge evolution and structural changes, enabling dynamic regrouping, cluster reformulation, and the emergence of new knowledge communities.

\subsection{Generative AI in Interactive Teaching}
Bridging the gap between raw document extraction algorithms and end-user educational platforms is the integration of Generative AI and Large Language Models (LLMs). The increasing accessibility and reasoning capabilities of Generative AI tools have led to widespread exploration in educational settings, fundamentally transforming traditional teaching and learning paradigms~\cite{AI_in_Edu}. Modern LLM agents have proven capable of autonomous reasoning, tool manipulation, and orchestrating complex academic tasks. 

Recent surveys highlight the versatility of LLM agents in automating multifaceted educational workflows, ranging from curriculum design to generating highly contextualized feedback~\cite{llm_in_edu}. These agents operate dynamically, leveraging multi-hop reasoning and continuous interaction to adapt to a student's cognitive state. A qualitative study involving educators demonstrated that the core impact of generative models spans three critical pillars: active teaching and learning, administrative automation, and personalized assessments~\cite{AI_in_Edu}. 

A prominent large-scale empirical example of this shift is the integration of AI within Harvard University's CS50 course. By deploying a suite of specialized AI software, the course approximated a 1:1 teacher-to-student ratio. Crucially, the system was designed with pedagogical boundaries, guiding students toward conceptual understanding rather than merely offering outright code solutions. This constrained autonomy proved highly effective in code explanation, style improvement, and administrative querying, demonstrating the practical viability of AI as a ubiquitous teaching assistant~\cite{ai_4_CS50}.

Despite these opportunities, deploying LLM agents in educational ecosystems introduces substantial systemic and ethical challenges. The probabilistic nature of LLMs frequently leads to hallucination, where the model generates factually incorrect or flawed information. Theoretical analyses of LLM representations reveal phenomena such as "semantic collapse," where the continuous manifold of the model's activations struggles to maintain stable, logically sound ontologies over continuous, long-term reasoning tasks~\cite{semantic_collapse}. 

To address this limitation, research in educational AI applications has shifted from static next-token prediction to autonomous agentic systems, aiming for frameworks that tightly navigate semantic boundaries in an academic context to ensure the reliability, alignment, and transparency of AI outputs. Architectures like the Adaptive Evaluation Multi-Agent (AEMA) framework have been proposed to address this gap. AEMA provides an auditable, multi-step evaluation mechanism operating under human oversight, ensuring that agents' internal decision-making processes, tool selections, and collaborative strategies remain verifiable and transparent~\cite{AEMA_agentic_eval}. Implementing such verifiable frameworks is essential to mitigate the risks of unconstrained autonomy and ensure that AI integration in education remains trustworthy and pedagogically effective.

\subsection{Educational Intelligent System}
Intelligent educational systems have transitioned from pure research into commercial applications. Many intelligent commercial platforms have emerged recently, alongside research institutions proposing full end-to-end pipelines beyond standard benchmark studies. 

In the research domain, systems such as~\cite{EduKG} (along with variants like K12EduKG and MOOC-KG) have been developed to uniformly model knowledge from textbooks and online courses. The~\cite{EduKGPipe} framework automates KG construction from unstructured PDF documents by extracting keywords via KeyBERT and performing semantic linking with LLMs. Furthermore, \texttt{MEduKG} pushes beyond text limitations by integrating multi-modal data into the graph, thereby enhancing the system's comprehensiveness and engagement for students.

In the commercial sector, popular platforms such as \texttt{Khanmigo} leverage generative LLMs to provide personalized tutoring and real-time feedback across extensive curriculum resources. In Vietnam, platforms such as \texttt{VioEdu} utilize manually curated knowledge trees to structure curricula, while \texttt{EdMicro}, specialized for enterprise training, employs skill matrices to map learning objectives and prerequisites. These systems represent distinct approaches to AI integration in education, each with unique strengths and limitations regarding scalability, personalization, and content structuring. Based on practical surveys and evaluation criteria, current educational systems exhibit diverse characteristics in AI adoption and knowledge structuring. A comparative analysis of these platforms is presented in Table~\ref{tab:edu_sys}.

\input{table/chap2/edu_sys.tex}

\subsection{Conclusion and Research Gaps}\label{sub:gaps}
Based on an extensive survey of the representative literature, including peer reviews and author self-evaluations, this report pinpoints the following core research gaps in the current landscape:

\paragraph*{{Limitations in Automation and Cost Optimization:}} Existing large-scale systems heavily depend on manually constructed knowledge trees or skill matrices. While course content is typically structured and ordered by domain expertise indexing, the lack of deep internal semantic modeling renders these systems less adaptive to individual student needs. This often results in effective high-level course recommendations but poor micro-level concept recommendations. Consequently, learning paths remain linear and rigid, with minimal real-time adaptation or automated content revision. This lack of automation restricts system scalability, as manual labeling becomes increasingly impractical as content repositories grow.
\paragraph*{{Lack of Connectivity and Transparency in Lesson Structures:}} Traditional text-mining and clustering methods often produce disconnected or fragmented knowledge groups. Existing learning platforms typically lack a standardized, mathematically verified network algorithm for modularizing study topics, resulting in disjointed learning paths that disrupt learners' cognitive flow.
\paragraph*{{Static Roadmap Reasoning and Gap Detection:}} While current Learning Management Systems (LMS) and AI tutors provide automated guidance, they rely primarily on traditional document retrieval or standard RAG, lacking the network structure and deterministic reasoning required to map prerequisites dynamically. The identification of critical "knowledge gaps" remains static and reactive, failing to automatically trace bottlenecks back to fundamental foundational concepts using network centrality metrics.

To date, frameworks such as~\cite{EduKGPipe} and its lineages have made significant attempts to address these issues through an end-to-end pipeline that fuses traditional text extraction with transformer-based weighting. Conversely, approaches like~\cite{DKG} propose what~\cite{EduKGPipe} lacks: an efficient mathematical community detection method that resolves the connectivity and transparency issues in lesson structures encountered by earlier pipelines. However, both methodologies exhibit distinct limitations. The former lacks deterministic knowledge mapping, causing knowledge disruption even within unfragmented graphs; the latter is less efficient in multi-step extraction, making it valuable for static knowledge modeling and analysis but suboptimal for real-time inference and recommendation. Finally, many theoretical proposals remain unadopted due to the absence of a comprehensive, production-grade platform capable of integrating these research advancements. 

